\hypertarget{c20-ranges}{%
\chapter{C20 RANGES}\label{c20-ranges}}

\hypertarget{c20-ranges-1}{%
\chapter{C++20 RANGES}\label{c20-ranges-1}}

\hypertarget{the-end-of-iterator-pairs}{%
\section{1. The End of Iterator pairs}\label{the-end-of-iterator-pairs}}

Ranges allow operating on a ``range'' (something with a begin and end)
directly, composing algorithms using pipe syntax
(\passthrough{\lstinline!|!}).

\begin{lstlisting}[language={C++}]
#include <ranges>
#include <vector>
#include <iostream>
#include <algorithm>

namespace views = std::views;
namespace ranges = std::ranges;

int main() {
    std::vector<int> nums = {1, 2, 3, 4, 5, 6};

    // Old Way
    // std::vector<int> temp;
    // std::copy_if(nums.begin(), nums.end(), std::back_inserter(temp), [](int i){ return i % 2 == 0; });
    // for(auto& x : temp) x *= x;

    // C++20 Ranges
    auto result = nums 
        | views::filter([](int i){ return i % 2 == 0; }) // Keep evens
        | views::transform([](int i){ return i * i; });  // Square them
        
    for (int i : result) {
        std::cout << i << " "; // 4 16 36
    }
}
\end{lstlisting}

\hypertarget{views-vs-actions}{%
\section{2. Views vs Actions}\label{views-vs-actions}}

\begin{itemize}
\tightlist
\item
  \textbf{Views:} Lazy. They adapt the iteration logic but don't own
  data. Computation happens \emph{during} iteration. (e.g.,
  \passthrough{\lstinline!views::filter!},
  \passthrough{\lstinline!views::reverse!}).
\item
  \textbf{Actions:} Eager. They modify the container immediately. (Not
  fully standardized in C++20, mostly views).
\end{itemize}

\hypertarget{projections}{%
\section{3. Projections}\label{projections}}

Algorithms now accept a ``projection'' to transform data before
comparison.

\begin{lstlisting}[language={C++}]
struct User { int id; std::string name; };
std::vector<User> users = {{2, "Bob"}, {1, "Alice"}};

// Sort by ID
ranges::sort(users, {}, &User::id);

// Sort by Name
ranges::sort(users, {}, &User::name);
\end{lstlisting}

\hypertarget{dangling-iterators-protection}{%
\section{4. Dangling Iterators
Protection}\label{dangling-iterators-protection}}

Ranges algorithms prevent returning iterators to temporary objects that
would die immediately.

\begin{lstlisting}[language={C++}]
auto get_vec() { return std::vector{1, 2, 3}; }

auto it = ranges::max_element(get_vec()); // Error! 
// Returns std::ranges::dangling instead of an invalid iterator.
\end{lstlisting}

\begin{center}\rule{0.5\linewidth}{0.5pt}\end{center}
